% ---------------------------------------------------------------------
\section{Descriptive Modeling}
% ---------------------------------------------------------------------

\subsection{Word Analysis \& Embeddings}

Text-based features were constructed through TF-IDF with dimensionality reduction, word embeddings, and topic modeling. For TF-IDF, each day's text was aggregated per platform into a single document. The vectorizer used unigrams and bigrams, a vocabulary cap of 500 features, logarithmic term frequency scaling (\texttt{sublinear\_tf=True}), and English stopword removal. Critically, the vectorizer was fitted exclusively on training data within each cross-validation fold to prevent vocabulary leakage from future documents (Lecture~2). The resulting sparse matrix was reduced to five dense components per platform via Singular Value Decomposition, yielding 15 TF-IDF features in total.

For word embeddings, a Word2Vec skip-gram model was trained from scratch on the news headline corpus, with separate models per political lean (Democratic, Republican, center), enabling cross-partisan semantic comparison \citep{mikolov2013}. The model used 100 dimensions, a context window of five, and 30 epochs. Pre-trained GloVe embeddings (\texttt{glove.6B.100d}) complemented this approach, covering 97.4\% of headline tokens and providing better generalization for low-frequency terms \citep{pennington2014}. Both spaces were visualized using t-SNE (perplexity=30) for local cluster structure and UMAP (\texttt{n\_neighbors=15}, cosine metric) for global relationships \citep{vandermaaten2008}.

Topic modeling was applied to Reddit posts using LDA \citep{blei2003}. After filtering (\texttt{no\_below=5}, \texttt{no\_above=0.5}), the vocabulary contained 12,302 words across 85,446 documents. A coherence sweep over $K=2$ to $10$ returned $K=2$ as optimal (\texttt{u\_mass}$=-3.54$), reflecting the thematic homogeneity of campaign text in which discourse reduces to a pro-Trump versus pro-Harris axis. The model was fitted with 15 passes and 400 EM iterations.

\subsection{Sentiment Analysis}

Five sentiment methods were applied across the three text sources to capture distinct affective dimensions of political discourse. The multi-method design allows the predictive model to select the most informative signal through regularization rather than requiring a single method to generalize across all contexts (Lectures~4--5).

VADER was applied to all platforms as the primary lexicon-based method. Its trigram-based valence shifting mechanism handles negations, intensifiers, capitalization, and punctuation, making it well-suited to social media text \citep{hutto2014}. Twitter-RoBERTa (\texttt{cardiffnlp/twitter-roberta-base-sentiment-latest}) complemented VADER on social media by processing full token sequences bidirectionally, capturing contextual dependencies that VADER's trigram window cannot resolve \citep{barbieri2022}. The pairwise correlation between their compound scores was 0.54, confirming partial but non-redundant overlap. For news headlines, FinBERT replaced VADER as the primary method given its domain-specific pre-training on financial text, which closely resembles the forward-looking and evaluative framing of campaign coverage. TextBlob provided a subjectivity dimension orthogonal to polarity (correlation with VADER: 0.43), and NRCLex extracted eight primary emotions from Plutchik's wheel to capture affect beyond positive/negative polarity (correlation with VADER: 0.34). Finally, SBERT (\texttt{all-MiniLM-L6-v2}) embedded 50 random sentences per day and source into 384-dimensional vectors, reduced to five principal components via PCA, yielding 15 dense semantic features with frozen weights.

All methods are summarized in the Appendix. Sentiment scores were aggregated daily and computed separately for TrumpBuzz and HarrisBuzz documents to produce sentiment gap features. A key descriptive finding is that the news sentiment gap correlates at $r=0.31$ with Polymarket prices, while raw news volume correlates near zero, suggesting that affective framing rather than volume carries market-relevant information.

\subsection{Network Analysis}

A directed mention graph was constructed from Bluesky data using NetworkX, with a weighted edge from account~A to~B for each @mention of~B by~A. After normalizing usernames by stripping the \texttt{.bsky.social} suffix, the full graph contains 1,467 nodes and 1,474 edges (density: 0.000685). The largest connected component (LCC) retains 928 nodes (63.3\%), with density 0.001317, diameter 15, mean shortest path 5.83 hops, and an average clustering coefficient of 0.0016 (Lecture~1). The low clustering coefficient confirms the broadcast-oriented nature of the platform: users mention others without sharing mutual audiences.

The degree distribution follows a power-law pattern on a log-log scale, confirming a scale-free topology in which a small number of hub nodes accumulate a disproportionate share of incoming mentions --- the most-mentioned account received 66 in-degree mentions. Hub nodes with high betweenness centrality act as information bridges whose activity may correlate with broader shifts in public discourse.

Echo chamber structure was quantified via a cross-cluster mention matrix partitioning users into TrumpBuzz, HarrisBuzz, and ElectionBuzz clusters. The diagonal dominance of this matrix operationalizes homophily: the principle that individuals preferentially associate with similar others \citep{hamilton2016}. Fewer than 6.8\% of LCC nodes bridge both political clusters. Weekly within-cluster ratios reveal monotonically increasing insularity approaching election day. The Reddit author overlap analysis corroborates this: the Jaccard similarity between \texttt{r/conservative} and \texttt{r/democrats} authors is only 0.019 (0.8\% cross-partisan participation), providing strong evidence for structurally enforced echo chambers \citep{hamilton2016}. Network-derived daily features including density, hub degree averages, community homophily, and echo chamber velocity were included in the predictive model.
